% ═══════════════════════════════════════════════════════════════════════════
% ARBEITSBLATT: Impuls & Impulserhaltung (DS2)
% Physik 10E Gymnasium | 26 BE
% ═══════════════════════════════════════════════════════════════════════════

\documentclass[11pt,a4paper]{article}

% ═══════════════════════════════════════════════════════════════════════════
% PAKETE
% ═══════════════════════════════════════════════════════════════════════════

\usepackage[utf8]{inputenc}
\usepackage[T1]{fontenc}
\usepackage[ngerman]{babel}
\usepackage{helvet}
\renewcommand{\familydefault}{\sfdefault}

\usepackage[margin=20mm,top=25mm,bottom=20mm]{geometry}
\usepackage{tikz}
\usepackage{xcolor}
\usepackage{amsmath}
\usepackage{amssymb}
\usepackage{tabularx}
\usepackage{enumitem}
\usepackage{fancyhdr}
\usepackage{lastpage}
\usepackage{tcolorbox}
\usepackage{qrcode}
\usepackage{hyperref}

\tcbuselibrary{skins,breakable}

% URL für QR-Code
\newcommand{\lpurl}{https://devbox42.github.io/sciencesim/physik/kl10/06-bewegungslehre-v2/LP-02-impuls.html}

% ═══════════════════════════════════════════════════════════════════════════
% FARBEN
% ═══════════════════════════════════════════════════════════════════════════

\definecolor{physik}{HTML}{2c5aa0}
\definecolor{physik-light}{HTML}{e8f0fa}
\definecolor{success}{HTML}{2a7a4b}
\definecolor{warning}{HTML}{b35c00}

% ═══════════════════════════════════════════════════════════════════════════
% KOPF-/FUSSZEILE
% ═══════════════════════════════════════════════════════════════════════════

\pagestyle{fancy}
\fancyhf{}
\renewcommand{\headrulewidth}{0.5pt}
\fancyhead[L]{\small\textcolor{physik}{\textbf{Physik 10E}} | Impuls \& Impulserhaltung}
\fancyhead[R]{\small Name: \underline{\hspace{5cm}}}
\fancyfoot[C]{\small Seite \thepage\ von \pageref{LastPage}}

% ═══════════════════════════════════════════════════════════════════════════
% BEFEHLE
% ═══════════════════════════════════════════════════════════════════════════

% Lücke
\newcommand{\luecke}[1]{\underline{\hspace{#1}}}

% Antwortzeilen
\newcommand{\antwortzeilen}[1]{%
    \par\vspace{3pt}%
    \foreach \i in {1,...,#1} {\noindent\rule{\linewidth}{0.4pt}\par\vspace{6pt}}%
}

% Kasten für Definitionen (nur Rahmen farbig, Inhalt weiß)
\newtcolorbox{kasten}[1][]{%
    colback=white,
    colframe=physik,
    colbacktitle=white,
    coltitle=black,
    fonttitle=\bfseries,
    title=#1,
    boxrule=1pt,
    arc=0pt,
    left=8pt,
    right=8pt,
    top=6pt,
    bottom=6pt
}

% Formelkasten
\newtcolorbox{formelbox}{%
    colback=white,
    colframe=physik,
    boxrule=2pt,
    arc=0pt,
    left=10pt,
    right=10pt,
    top=10pt,
    bottom=10pt
}

% AFB-Marker (dezent in Grau)
\newcommand{\afb}[1]{\hfill\textcolor{gray}{\scriptsize AFB #1}}

% Punkte am Rand
\newcommand{\punkte}[1]{\hfill\textcolor{gray}{\small(#1 BE)}}

% ═══════════════════════════════════════════════════════════════════════════
% DOKUMENT
% ═══════════════════════════════════════════════════════════════════════════

\begin{document}

% ───────────────────────────────────────────────────────────────────────────
% TITEL
% ───────────────────────────────────────────────────────────────────────────

\noindent
\begin{minipage}[c]{0.82\textwidth}
    \begin{center}
        {\Large\bfseries\textcolor{physik}{Arbeitsblatt: Impuls \& Impulserhaltung}}\\[5pt]
        {\small DS2 | Lernpfad: LP-02-impuls}
    \end{center}
\end{minipage}%
\hfill
\begin{minipage}[c]{0.15\textwidth}
    \raggedleft
    \href{\lpurl}{\qrcode[height=15mm]{\lpurl}}\\[2pt]
    {\tiny\textcolor{gray}{Lernpfad}}
\end{minipage}

\vspace{5mm}

% ═══════════════════════════════════════════════════════════════════════════
% KASTEN K1: DEFINITION
% ═══════════════════════════════════════════════════════════════════════════

\begin{kasten}[K1: Was ist Impuls? \afb{I}]

Vervollständige den Lückentext:

\vspace{8pt}

Der \textbf{Impuls} beschreibt, wie schwer es ist, einen bewegten Körper zu

\luecke{3cm}. Je größer der Impuls, desto mehr \glqq\luecke{2cm}\grqq{} hat der Körper.

\vspace{8pt}

Der Impuls hängt von zwei Größen ab: der \luecke{2.5cm} und der \luecke{3.5cm}.

\end{kasten}

\vspace{5mm}

% ═══════════════════════════════════════════════════════════════════════════
% KASTEN K2: FORMEL
% ═══════════════════════════════════════════════════════════════════════════

\begin{kasten}[K2: Formel und Einheit \afb{I}]

\begin{center}
\begin{formelbox}
\vspace{5pt}
{\Large $p = $ \luecke{1.5cm} $\cdot$ \luecke{1.5cm}}

\vspace{10pt}

{\small Einheit: $[p] = $ \luecke{3cm}}
\vspace{5pt}
\end{formelbox}
\end{center}

\vspace{3pt}

\begin{tabular}{@{}l@{\hspace{5pt}}l@{\hspace{15pt}}l@{\hspace{5pt}}l@{\hspace{15pt}}l@{\hspace{5pt}}l@{}}
$p$ = & Impuls & $m$ = & \luecke{2cm} & $v$ = & \luecke{3cm}
\end{tabular}

\end{kasten}

\vspace{5mm}

% ═══════════════════════════════════════════════════════════════════════════
% ÜBUNG Ü1: IMPULS BERECHNEN
% ═══════════════════════════════════════════════════════════════════════════

\textbf{Ü1: Impuls berechnen} \afb{I/II}

\vspace{3mm}

\begin{enumerate}[label=\alph*), leftmargin=*, itemsep=12pt]

\item Ein \textbf{Fußball} (m = 0,4 kg) wird mit v = 25 m/s geschossen.

\vspace{3pt}
$p = m \cdot v = $ \luecke{5cm} $= $ \luecke{3cm}

\item Ein \textbf{Sprinter} (m = 80 kg) rennt mit v = 10 m/s.

\vspace{3pt}
$p = $ \luecke{7cm}

\item Ein \textbf{Auto} (m = 1200 kg) fährt mit v = 30 m/s ($\approx$ 108 km/h).

\vspace{3pt}
$p = $ \luecke{7cm}

\item Ein \textbf{ICE} (m = 400.000 kg) fährt mit v = 80 m/s ($\approx$ 288 km/h).

\vspace{3pt}
$p = $ \luecke{7cm}

\vspace{5pt}
\textit{Warum ist ein Zug gefährlicher als ein Auto, obwohl er oft langsamer fährt?}

\vspace{3pt}
\antwortzeilen{2}

\end{enumerate}

\newpage

% ═══════════════════════════════════════════════════════════════════════════
% ÜBUNG Ü2: VERGLEICHE
% ═══════════════════════════════════════════════════════════════════════════

\textbf{Ü2: Wer hat mehr Wucht?} \afb{II}

\vspace{3mm}

\begin{enumerate}[label=\alph*), leftmargin=*, itemsep=10pt]

\item Berechne und vergleiche:

\vspace{5pt}
\begin{tabular}{|l|c|c|c|}
\hline
\textbf{Objekt} & \textbf{Masse} & \textbf{Geschwindigkeit} & \textbf{Impuls} \\
\hline
Usain Bolt & 94 kg & 12 m/s & \luecke{2.5cm} \\
\hline
Sumo-Ringer & 180 kg & 4 m/s & \luecke{2.5cm} \\
\hline
\end{tabular}

\vspace{5pt}
Wer hat mehr Impuls? \luecke{3cm} \hfill \textit{Begründung:} \luecke{5cm}

\vspace{8pt}

\item Berechne DEINEN Impuls beim Gehen (v $\approx$ 1,5 m/s):

\vspace{3pt}
Meine geschätzte Masse: \luecke{2cm} kg \hfill Mein Impuls: $p = $ \luecke{4cm}

\vspace{8pt}

\item \textbf{Transfer:} Erkläre, warum Airbags bei einem Unfall helfen, obwohl sie den Impuls nicht ändern können. \textit{(Tipp: Denke an ZEIT)} \afb{III}

\vspace{3pt}
\antwortzeilen{3}

\end{enumerate}

\vspace{5mm}

% ═══════════════════════════════════════════════════════════════════════════
% KASTEN K3: MERKSATZ
% ═══════════════════════════════════════════════════════════════════════════

\begin{kasten}[K3: Merksatz Teil A \afb{I}]

Vervollständige:

\vspace{8pt}

Der \textbf{Impuls} $p = m \cdot v$ beschreibt die \glqq\luecke{2cm}\grqq{} eines bewegten Körpers.

\vspace{8pt}

Je größer die \luecke{2.5cm} \textbf{oder} die \luecke{3cm}, desto größer ist der Impuls.

\end{kasten}

\newpage

% ═══════════════════════════════════════════════════════════════════════════
% TEIL B: IMPULSERHALTUNG
% ═══════════════════════════════════════════════════════════════════════════

{\large\bfseries\textcolor{physik}{Teil B: Impulserhaltung}}

\vspace{5mm}

% ═══════════════════════════════════════════════════════════════════════════
% KASTEN K4: ABGESCHLOSSENES SYSTEM
% ═══════════════════════════════════════════════════════════════════════════

\begin{kasten}[K4: Abgeschlossenes System \afb{I}]

Ein \textbf{abgeschlossenes System} ist ein Bereich, in dem \luecke{4cm} wirken.

\vspace{5pt}

Beispiele: \luecke{6cm}

\end{kasten}

\vspace{5mm}

% ═══════════════════════════════════════════════════════════════════════════
% KASTEN K5: IMPULSERHALTUNGSSATZ
% ═══════════════════════════════════════════════════════════════════════════

\begin{kasten}[K5: Impulserhaltungssatz \afb{I}]

\begin{center}
\begin{formelbox}
\vspace{5pt}
{\Large $p_{\text{vorher}} = $ \luecke{3cm}}

\vspace{10pt}

{\normalsize In einem abgeschlossenen System bleibt der \textbf{Gesamtimpuls} \luecke{2.5cm}.}
\vspace{5pt}
\end{formelbox}
\end{center}

\vspace{3pt}

Ausführlich: $m_1 \cdot v_1 + m_2 \cdot v_2 = m_1 \cdot v_1' + m_2 \cdot v_2'$

\end{kasten}

\vspace{5mm}

% ═══════════════════════════════════════════════════════════════════════════
% KASTEN K6: VERGLEICH STOSSARTEN
% ═══════════════════════════════════════════════════════════════════════════

\begin{kasten}[K6: Elastischer vs. Unelastischer Stoß \afb{II}]

\vspace{3pt}

\begin{tabular}{|p{4cm}|p{5cm}|p{5cm}|}
\hline
& \textbf{Elastischer Stoß} & \textbf{Unelastischer Stoß} \\
\hline
\textbf{Was passiert?} & Körper & Körper \\[0.8em]
\hline
\textbf{Energie} & bleibt erhalten & geht verloren (Wärme) \\[0.8em]
\hline
\textbf{Impuls} & & \\[0.8em]
\hline
\textbf{Beispiel} & & \\[0.8em]
\hline
\end{tabular}

\end{kasten}

\newpage

% ═══════════════════════════════════════════════════════════════════════════
% ÜBUNG Ü3: ELASTISCHER STOSS
% ═══════════════════════════════════════════════════════════════════════════

\textbf{Ü3: Elastischer Stoß} \afb{II}

\vspace{3mm}

Zwei gleich schwere Wagen auf einer Luftkissenbahn:
\begin{itemize}[leftmargin=*, itemsep=2pt]
    \item Wagen 1: $m_1 = 2$ kg, $v_1 = 3$ m/s (fährt nach rechts)
    \item Wagen 2: $m_2 = 2$ kg, $v_2 = 0$ m/s (steht still)
\end{itemize}

Nach dem \textbf{elastischen} Stoß steht Wagen 1 still ($v_1' = 0$).

\vspace{5pt}

a) Gesamtimpuls vorher: $p_{\text{ges}} = m_1 \cdot v_1 + m_2 \cdot v_2 = $ \luecke{4cm}

\vspace{5pt}

b) Nach dem Stoß: $p_{\text{ges}} = m_2 \cdot v_2'$ \hfill $\Rightarrow$ $v_2' = $ \luecke{3cm}

\vspace{8mm}

% ═══════════════════════════════════════════════════════════════════════════
% ÜBUNG Ü4: UNELASTISCHER STOSS
% ═══════════════════════════════════════════════════════════════════════════

\textbf{Ü4: Unelastischer Stoß (Güterwagen)} \afb{II}

\vspace{3mm}

Ein Güterwagen ($m_1 = 20.000$ kg, $v_1 = 5$ m/s) rollt auf einen stehenden Wagen ($m_2 = 30.000$ kg).

Nach dem Kuppeln bewegen sich beide gemeinsam weiter.

\vspace{5pt}

a) Gesamtimpuls vorher: $p_{\text{ges}} = m_1 \cdot v_1 = $ \luecke{4cm}

\vspace{5pt}

b) Gemeinsame Geschwindigkeit: $v' = \dfrac{p_{\text{ges}}}{m_1 + m_2} = $ \luecke{4cm}

\vspace{8mm}

% ═══════════════════════════════════════════════════════════════════════════
% ÜBUNG Ü5: TRANSFER
% ═══════════════════════════════════════════════════════════════════════════

\textbf{Ü5: Transfer – Sicherheitstechnik} \afb{III}

\vspace{3mm}

Erkläre: Warum ist ein Aufprall gegen eine \textbf{Knautschzone} weniger gefährlich als gegen eine starre Wand? \textit{(Tipp: Denke an Impulsübertragung und ZEIT)}

\vspace{3pt}
\antwortzeilen{4}

\vspace{8mm}

% ═══════════════════════════════════════════════════════════════════════════
% KASTEN K7: ZUSAMMENFASSUNG
% ═══════════════════════════════════════════════════════════════════════════

\begin{kasten}[K7: Zusammenfassung \afb{I}]

\textbf{Teil A:} Der Impuls $p = $ \luecke{2cm} beschreibt die \glqq Wucht\grqq{} eines bewegten Körpers.

\vspace{5pt}

\textbf{Teil B:} In einem abgeschlossenen System gilt: $p_{\text{vorher}} = $ \luecke{2.5cm}

\vspace{5pt}

\begin{tabular}{@{}ll@{}}
Elastischer Stoß: & Körper prallen \luecke{1.5cm} \\[0.3em]
Unelastischer Stoß: & Körper kleben \luecke{1.5cm}
\end{tabular}

\end{kasten}

\end{document}
